\section{Consequences for Robin's inequality}
\label{s-robin-margin}

The two preceding results, together with
\(-g_{\integers}(m)=\Igap(m)-g_{\reals}(m)\), give the integer margin.

\begin{theorem}[Fixed-slice Robin margin]
\label{thm-rh-margin}
Under the Riemann hypothesis,
\begin{equation}
\label{e-rh-margin}
    -\sqrt{p_m}\log p_m\,g_{\integers}(m)
    =
    2\sqrt2-2+\mathcal H(\log p_m)+o(1).
\end{equation}
Consequently,
\[
    g_{\reals}(m)>0>g_{\integers}(m)
\]
for every sufficiently large \(m\).

If the Riemann hypothesis is false, then
\[
\begin{aligned}
    \limsup_{m\to\infty}
    -\sqrt{p_m}\log p_m\,g_{\integers}(m)
    &=+\infty,\\
    \liminf_{m\to\infty}
    -\sqrt{p_m}\log p_m\,g_{\integers}(m)
    &=-\infty.
\end{aligned}
\]
\end{theorem}

\begin{proof}
Under the Riemann hypothesis, subtract
Theorem~\ref{thm-rh-relaxed} from the gap formula in
Theorem~\ref{thm-integrality-gap}. This gives \eqref{e-rh-margin}.
The relaxed value is eventually positive by
Theorem~\ref{thm-rh-relaxed}. The integer value is eventually negative
because
\[
    2\sqrt2-2+\mathcal H(\log p_m)
    \geq
    2\sqrt2-2-\tau
    >
    0.
\]

If the hypothesis is false, Theorem~\ref{thm-off-rh-relaxed} gives
opposite-sign excursions of size \(p_m^{-b}\) for every admissible
\(b<1/2\). The integrality gap is
\[
    O\left(\frac1{\sqrt{p_m}\log p_m}\right)
    =
    o(p_m^{-b}).
\]
It cannot suppress either excursion. Since
\(-g_{\integers}=\Igap-g_{\reals}\), the margin has the claimed
opposite-sign divergence.
\end{proof}

\begin{corollary}[Eventual fixed-\texorpdfstring{\(\omega\)}{omega}
criterion]
\label{cor-fixed-omega-equivalence}
The Riemann hypothesis is equivalent to the existence of an \(M\) such
that Robin's inequality holds for every \(n\geq5041\) with
\(\omega(n)=m\), for every \(m\geq M\).
\end{corollary}

\begin{proof}
If the Riemann hypothesis holds, the theorem gives
\(g_{\integers}(m)<0\) for all sufficiently large \(m\), so every
integer in each such slice satisfies Robin's inequality.

If it fails, the theorem gives infinitely many \(m\) with
\(g_{\integers}(m)>0\). Proposition~\ref{prop-support} shows that the
maximum is attained, so each such slice contains a violation.
\end{proof}

Formula \eqref{e-rh-margin} separates the sign quantitatively. The
analytic term contributes \(2-\mathcal H\), while the rounding
correction contributes \(2\sqrt2\). The strict inequality
\(2\sqrt2>2+\tau\) is why Robin's inequality survives on the integer
slice but fails for its real-exponent relaxation.
